% Reviewed translation: PDF pages 212--218 / printed pages 198--204.
% The original scan is authoritative; mathematical tokens were checked against it.
\MathBookSourcePage{212}
\section{二维 Stokes 问题的 "流函数--涡量--压力" 方法}
\label{sec:incompressible-sfvp}

二维 Stokes 方程组中广泛使用的 "流函数--涡量--压力" 表述
建立在两个恒等式之上:
\[
\operatorname{curl}(\operatorname{curl}\mathbf u)
=-\Delta\mathbf u+\operatorname{grad}(\operatorname{div}\mathbf u),
\qquad
\operatorname{curl}(\operatorname{curl}\phi)=-\Delta\phi.
\]
因此 Stokes 问题
\[
\left.
\begin{aligned}
-\nu\Delta\mathbf u+\operatorname{grad}p&=\mathbf f,\\
\operatorname{div}\mathbf u&=0
\end{aligned}
\right\}\quad\text{in }\Omega,
\qquad
\mathbf u|_\Gamma=\mathbf0,
\]
既可表示为
\[
\left.
\begin{aligned}
\operatorname{curl}\mathbf u&=\omega,
&\operatorname{div}\mathbf u&=0,\\
\nu\operatorname{curl}\omega+\operatorname{grad}p&=\mathbf f
\end{aligned}
\right\}\quad\text{in }\Omega,
\qquad
\mathbf u|_\Gamma=\mathbf0,
\]
也可表示为
\[
\left.
\begin{aligned}
\mathbf u&=\operatorname{curl}\psi,
&\omega&=-\Delta\psi,\\
\nu\operatorname{curl}\omega+\operatorname{grad}p&=\mathbf f
\end{aligned}
\right\}\quad\text{in }\Omega,
\qquad
\operatorname{curl}\psi|_\Gamma=\mathbf0.
\]
本节和下一节讨论的混合方法由这两个表述导出.

\subsection{一个混合表述}
\label{subsec:incompressible-sfvp-mixed-formulation}

令 $\Omega$ 为 $\mathbb R^2$ 中的有界区域,
其 Lipschitz 连续边界 $\Gamma$ 的连通分支如
Figure~\ref{fig:sec3-multiply-connected-domain} 所示,
记为 $\Gamma_i$, $0\leq i\leq p$.
给定 $\mathbf f\in L^r(\Omega)^2$, $r>1$,
回顾 $\Omega$ 中 Stokes 方程的经典表述:

求 $(\mathbf u,p)\in H_0^1(\Omega)^2\times L_0^2(\Omega)$, 使
\begin{equation}\label{eq:incompressible-sfvp-stokes}
\left\{
\begin{aligned}
\nu(\operatorname{grad}\mathbf u,\operatorname{grad}\mathbf v)
-(p,\operatorname{div}\mathbf v)
&=(\mathbf f,\mathbf v)
&&\forall\mathbf v\in H_0^1(\Omega)^2,\\
(q,\operatorname{div}\mathbf u)&=0
&&\forall q\in L_0^2(\Omega).
\end{aligned}
\right.
\end{equation}

\long\gdef\ChapterThreeSectionTwoDeferredContinuation{%
\MathBookSourcePage{215}
\begin{Theorem}\label{thm:incompressible-stream-vorticity-wellposed}
问题 \eqref{eq:incompressible-stream-vorticity-pair}
在 $\Psi\times L^2(\Omega)$ 中有唯一解
$(\psi,\omega=-\Delta\psi)$.
并且存在唯一的 $p\in L_0^2(\Omega)$,
使 $(\mathbf u=\operatorname{curl}\psi,\omega,p)$
是问题
\eqref{eq:incompressible-velocity-vorticity-pressure} 的解.
\end{Theorem}

然而, 这些问题都不能完全满足我们的目的,
因为其内逼近要求构造有限维子空间, 其中要么速度严格无散并具有
$H^1$ 正则性, 要么与之等价的流函数具有 $H^2$ 正则性.
实际中这非常不理想.
为解决这一困难, 我们放松测试速度 $\mathbf v$
或流函数 $\phi$ 的正则性.

先考虑问题 \eqref{eq:incompressible-sfvp-four-field}.
双线性形式 $a(\cdot,\cdot)$ 可以保持不变,
但包含 $\operatorname{curl}\mathbf v$ 和
$\operatorname{div}\mathbf v$ 的形式 $b(\cdot,\cdot)$ 必须修改.
对 $\mathbf v\in H_0^1(\Omega)^2$, 有
\[
\left.
\begin{aligned}
(\operatorname{curl}\mathbf v,\mu)
&=(\mathbf v,\operatorname{curl}\mu),\\
(q,\operatorname{div}\mathbf v)
&=-(\operatorname{grad}q,\mathbf v)
\end{aligned}
\right\}
\qquad\forall\mu,q\in H^1(\Omega).
\]
因此, 若提高 $\mu,q$ 的正则性, 就可以降低 $\mathbf v$ 的正则性.
最直接的选择是
$\mathbf v\in L^2(\Omega)^2$ 且 $\mu,q\in H^1(\Omega)$;
对于 Stokes 问题, 这已经相当合适.
但考虑到后续对 Navier--Stokes 方程的应用,
最好允许 $\mu,q$ 属于更广的一类空间.

为此, 固定实数 $s\geq2$, 并用
\[
\frac1r+\frac1s=1
\]
定义 $r\in\mathbb R$. 取
\begin{equation}\label{eq:incompressible-weak-four-field-spaces}
\widetilde X=L^s(\Omega)^2\times L^2(\Omega),
\qquad
\widetilde M=W^{1,r}(\Omega)
\times[W^{1,r}(\Omega)\cap L_0^2(\Omega)],
\end{equation}
相应范数为
\[
\|(\mathbf v,\theta)\|_{\widetilde X}
=\|\mathbf v\|_{0,s,\Omega}+\|\theta\|_{0,\Omega},
\qquad
\|(\mu,q)\|_{\widetilde M}
=\|\mu\|_{1,r,\Omega}+\|q\|_{1,r,\Omega},
\]
并定义双线性形式
\begin{equation}\label{eq:incompressible-weak-four-field-forms}
\left\{
\begin{aligned}
\widetilde a((\mathbf u,\omega),(\mathbf v,\theta))
&=\nu(\omega,\theta),\\
\widetilde b((\mathbf v,\theta),(\mu,q))
&=(\mathbf v,\operatorname{curl}\mu)
-(\theta,\mu)+(\operatorname{grad}q,\mathbf v),
\end{aligned}
\right.
\end{equation}
其中 $(\mathbf u,\omega),(\mathbf v,\theta)\in\widetilde X$,
$(\mu,q)\in\widetilde M$.

\MathBookSourcePage{216}
现在假设
\eqref{eq:incompressible-sfvp-stokes} 的右端项
$\mathbf f\in L^r(\Omega)^2$, 并考虑问题 $(\widetilde Q)$:

求 $(\mathbf u,\omega)\in\widetilde X$ 和
$(\lambda,p)\in\widetilde M$, 使
\begin{equation}\label{eq:incompressible-weak-four-field}
\left\{
\begin{aligned}
\nu(\omega,\theta)
+(\mathbf v,\operatorname{curl}\lambda)
-(\theta,\lambda)+(\operatorname{grad}p,\mathbf v)
&=(\mathbf f,\mathbf v)
&&\forall(\mathbf v,\theta)\in\widetilde X,\\
(\mathbf u,\operatorname{curl}\mu)
-(\omega,\mu)+(\operatorname{grad}q,\mathbf u)
&=0
&&\forall(\mu,q)\in\widetilde M.
\end{aligned}
\right.
\end{equation}
在继续之前, 先考察问题 $(\widetilde Q)$ 相联系的空间
\begin{equation}\label{eq:incompressible-weak-four-field-kernel}
\widetilde V
=\{(\mathbf v,\theta)\in\widetilde X;\
(\mathbf v,\operatorname{curl}\mu)-(\theta,\mu)
+(\operatorname{grad}q,\mathbf v)=0
\quad\forall(\mu,q)\in\widetilde M\}.
\end{equation}
在 $\widetilde X$ 上引入半范数
\begin{equation}\label{eq:incompressible-vorticity-seminorm}
|(\mathbf v,\theta)|=\|\theta\|_{0,\Omega}
\qquad\forall(\mathbf v,\theta)\in\widetilde X.
\end{equation}

\begin{Lemma}\label{lem:incompressible-weak-kernel-coincidence}
空间 $V$ 与 $\widetilde V$ 在代数和拓扑意义下相同.
此外, 半范数
\eqref{eq:incompressible-vorticity-seminorm}
在 $V$ 上是一个与 $\|\cdot\|_{\widetilde X}$ 等价的范数:
\[
|(\mathbf v,\theta)|
\simeq\|(\mathbf v,\theta)\|_{\widetilde X}
\qquad\forall(\mathbf v,\theta)\in V.
\]
\end{Lemma}

\begin{Proof}
设 $(\mathbf v,\theta)\in\widetilde V$.
直接应用 Green 公式依次得到
\[
\operatorname{div}\mathbf v=0\quad\text{in }\Omega,
\qquad
\mathbf v\cdot\mathbf n=0\quad\text{on }\Gamma,
\]
以及
\[
\theta=\operatorname{curl}\mathbf v\quad\text{in }\Omega,
\qquad
\mathbf v\cdot\boldsymbol\tau=0\quad\text{on }\Gamma.
\]
因而
$\mathbf v\in H_0(\operatorname{div};\Omega)
\cap H_0(\operatorname{curl};\Omega)$,
并由 Corollary~\ref{cor:sec2-local-h1-hdiv-hcurl} 得
\[
\mathbf v\in H_0^1(\Omega)^2,
\]
且
\begin{equation}\label{eq:incompressible-velocity-by-vorticity}
|\mathbf v|_{1,\Omega}
\leq C\|\theta\|_{0,\Omega}.
\end{equation}
所以 $\widetilde V\subset V$, 反向包含显然成立.
范数等价性由更强的结果
\eqref{eq:incompressible-velocity-by-vorticity}
和 Sobolev 嵌入 Theorem 立即得到.
\end{Proof}

现在可以证明问题 $(Q)$ 和 $(\widetilde Q)$ 等价.

\begin{Theorem}\label{thm:incompressible-weak-four-field-equivalence}
假设 Stokes 问题
\eqref{eq:incompressible-sfvp-stokes} 的解 $(\mathbf u,p)$ 满足
\begin{equation}\label{eq:incompressible-stokes-weak-regularity}
\operatorname{curl}\mathbf u\in W^{1,r}(\Omega),
\qquad
p\in W^{1,r}(\Omega)\cap L_0^2(\Omega).
\end{equation}
则问题 $(\widetilde Q)$ 有唯一解
\[
\mathbf u,\qquad
\omega=\operatorname{curl}\mathbf u,\qquad
\lambda=\nu\omega,\qquad
p.
\]
\end{Theorem}

\begin{Proof}
这是 Theorem~\ref{thm:incompressible-weak-mixed-equivalence}
和 Lemma~\ref{lem:incompressible-weak-kernel-coincidence}
的直接应用.
由 Sobolev 嵌入 Theorem 和相应的嵌入结果,
$X$ 在 $\widetilde X$ 中稠密且连续,
$\widetilde M$ 在 $M$ 中稠密且连续.
等式
\eqref{eq:incompressible-extended-a}
与 \eqref{eq:incompressible-extended-b} 显然成立.
椭圆性条件由
\eqref{eq:incompressible-velocity-by-vorticity} 得到,
而原有的两个假设此前已验证.
结合
\eqref{eq:incompressible-stokes-weak-regularity}
与 Lemma~\ref{lem:incompressible-weak-kernel-coincidence},
Theorem~\ref{thm:incompressible-weak-mixed-equivalence}
的全部条件均成立.
\end{Proof}

\MathBookSourcePage{217}
\begin{Remark}\label{rem:incompressible-weak-regularity-sufficient}
若 $\Gamma$ 属于 $\mathcal C^2$ 类,
或 $\Omega$ 是凸多边形,
则相应的 Stokes 正则性结果表明,
只要 $\mathbf f\in L^r(\Omega)^2$, 就有
$\mathbf u\in W^{2,r}(\Omega)^2$ 和
$p\in W^{1,r}(\Omega)$.
因此, 对 $\mathbf f$ 和 $\Omega$ 的这一温和正则性假设足以保证
问题 $(\widetilde Q)$ 有唯一解.
\end{Remark}

关于流函数, 回顾
Corollary~\ref{cor:sec3-h-stream-functions}:
$\operatorname{curl}$ 是从
\[
\Phi=\{\chi\in H^1(\Omega);\
\chi|_{\Gamma_0}=0,\
\chi|_{\Gamma_i}=c_i\text{ 为常数},\
1\leq i\leq p\}
\]
到
\[
H=\{\mathbf v\in L^2(\Omega)^2;\
\operatorname{div}\mathbf v=0,\
\mathbf v\cdot\mathbf n|_\Gamma=0\}
\]
上的同构.
这一结果立即蕴含,
$\operatorname{curl}$ 也是从
\begin{equation}\label{eq:incompressible-stream-space-phi-s}
\Phi_s=\Phi\cap W^{1,s}(\Omega)
\end{equation}
到
\[
H\cap L^s(\Omega)^2
=\{\mathbf v\in L^s(\Omega)^2;\
\operatorname{div}\mathbf v=0,\
\mathbf v\cdot\mathbf n|_\Gamma=0\}
\]
上的同构.
因此, 若 $((\mathbf u,\omega),(\lambda,p))$
是 \eqref{eq:incompressible-weak-four-field} 的解,
则一方面由
Lemma~\ref{lem:incompressible-weak-kernel-coincidence}
有 $(\mathbf u,\omega)\in V$;
另一方面, 在其中对每个 $\phi\in\Phi_s$
取 $\mathbf v=\operatorname{curl}\phi$, 可得
$\mathbf u=\operatorname{curl}\psi$, 其中
$((\psi,\omega),\lambda)$ 是下列问题的解:

求 $(\psi,\omega)\in\Phi_s\times L^2(\Omega)$ 和
$\lambda\in W^{1,r}(\Omega)$, 使
\begin{equation}\label{eq:incompressible-weak-stream-three-field}
\left\{
\begin{aligned}
\nu(\omega,\theta)-(\lambda,\theta)
+(\operatorname{curl}\phi,\operatorname{curl}\lambda)
&=(\mathbf f,\operatorname{curl}\phi)
&&\forall(\phi,\theta)\in\Phi_s\times L^2(\Omega),\\
(\operatorname{curl}\psi,\operatorname{curl}\mu)
-(\omega,\mu)&=0
&&\forall\mu\in W^{1,r}(\Omega).
\end{aligned}
\right.
\end{equation}
反之, 容易检验
\eqref{eq:incompressible-weak-stream-three-field}
至多有一个解. 因而得到如下两个重要结论.

\begin{Corollary}\label{cor:incompressible-weak-stream-unique}
在 Theorem~\ref{thm:incompressible-weak-four-field-equivalence}
的假设下, 问题
\eqref{eq:incompressible-weak-stream-three-field} 有唯一解
\[
\psi,\qquad
\omega=-\Delta\psi,\qquad
\lambda=\nu\omega,
\]
其中
$\operatorname{curl}\psi=\mathbf u$,
$\mathbf u$ 是
\eqref{eq:incompressible-sfvp-stokes} 的解.
\end{Corollary}

\begin{Theorem}\label{thm:incompressible-decoupled-stream-pressure}
在 Theorem~\ref{thm:incompressible-weak-four-field-equivalence}
的假设下, 问题
\eqref{eq:incompressible-weak-stream-three-field}
等价于:

求 $\psi\in\Phi_s$ 和 $\omega\in W^{1,r}(\Omega)$, 使
\begin{equation}\label{eq:incompressible-decoupled-stream-vorticity}
\left\{
\begin{aligned}
\nu(\operatorname{curl}\omega,\operatorname{curl}\phi)
&=(\mathbf f,\operatorname{curl}\phi)
&&\forall\phi\in\Phi_s,\\
(\operatorname{curl}\psi,\operatorname{curl}\mu)
&=(\omega,\mu)
&&\forall\mu\in W^{1,r}(\Omega).
\end{aligned}
\right.
\end{equation}

\MathBookSourcePage{218}
此外, 压力 $p$ 是下列问题的唯一解:

求 $p\in W^{1,r}(\Omega)\cap L_0^2(\Omega)$, 使
\begin{equation}\label{eq:incompressible-decoupled-pressure}
(\operatorname{grad}p,\operatorname{grad}q)
=(\mathbf f-\nu\operatorname{curl}\omega,
\operatorname{grad}q)
\qquad\forall q\in W^{1,s}(\Omega).
\end{equation}
\end{Theorem}

\begin{Proof}
在 \eqref{eq:incompressible-weak-four-field} 中,
对每个 $q\in W^{1,s}(\Omega)$ 取
$\mathbf v=\operatorname{grad}q$,
并利用 $\lambda=\nu\omega$, 即得到
\eqref{eq:incompressible-decoupled-pressure}.
\end{Proof}

后面各小节讨论的逼近主要建立在
\eqref{eq:incompressible-decoupled-stream-vorticity}
和 \eqref{eq:incompressible-decoupled-pressure} 上.
需要指出,
\eqref{eq:incompressible-decoupled-stream-vorticity}
可解释为 Laplace 算子的两个 Dirichlet 问题.
为简单起见, 考虑 $\Omega$ 单连通($p=0$)的情形.
令 $\omega_0$ 表示 $\omega$ 在 $\Gamma$ 上的迹,
则 \eqref{eq:incompressible-decoupled-stream-vorticity}
可解释为
\[
-\nu\Delta\omega=\operatorname{curl}\mathbf f
\quad\text{in }\Omega,
\qquad
\omega|_\Gamma=\omega_0,
\]
\[
-\Delta\psi=\omega
\quad\text{in }\Omega,
\qquad
\psi|_\Gamma=0.
\]
类似地,
\eqref{eq:incompressible-decoupled-pressure}
是 Laplace 算子的一个 Neumann 问题.
虽然
\eqref{eq:incompressible-decoupled-stream-vorticity}
中的两个问题相互耦合(因为 $\omega_0$ 未知),
但它与
\eqref{eq:incompressible-decoupled-pressure}
完全解耦: 先求解前者, 再求解后者.
此外, 这两个表述非常适合构造具有下列性质的逼近:

\textup{1)} 速度严格无散;

\textup{2)} 压力连续.
}
\MathBookSourcePage{213}
由于 $\operatorname{div}\mathbf u=0$, 已知
\begin{equation}\label{eq:incompressible-grad-curl-identity}
(\operatorname{grad}\mathbf u,\operatorname{grad}\mathbf v)
=(\operatorname{curl}\mathbf u,\operatorname{curl}\mathbf v)
\qquad\forall\mathbf v\in H_0^1(\Omega)^2.
\end{equation}
引入涡量 $\omega=\operatorname{curl}\mathbf u$,
方程 \eqref{eq:incompressible-sfvp-stokes} 等价于:

求 $(\mathbf u,\omega,p)
\in H_0^1(\Omega)^2\times L^2(\Omega)\times L_0^2(\Omega)$, 使
\begin{equation}\label{eq:incompressible-velocity-vorticity-pressure}
\left\{
\begin{aligned}
\nu(\omega,\operatorname{curl}\mathbf v)
-(p,\operatorname{div}\mathbf v)
&=(\mathbf f,\mathbf v)
&&\forall\mathbf v\in H_0^1(\Omega)^2,\\
(\operatorname{curl}\mathbf u-\omega,\mu)
-(q,\operatorname{div}\mathbf u)
&=0
&&\forall\mu\in L^2(\Omega),\quad
\forall q\in L_0^2(\Omega).
\end{aligned}
\right.
\end{equation}
把问题 \eqref{eq:incompressible-velocity-vorticity-pressure}
置于 Subsection~\ref{subsec:abstract-general-result} 的框架中.
置
\[
X=H_0^1(\Omega)^2\times L^2(\Omega),
\qquad
\|(\mathbf v,\theta)\|_X
=\bigl(|\mathbf v|_{1,\Omega}^2+\|\theta\|_{0,\Omega}^2\bigr)^{1/2},
\]
以及
\[
M=L^2(\Omega)\times L_0^2(\Omega),
\qquad
\|(\mu,q)\|_M
=\bigl(\|\mu\|_{0,\Omega}^2+\|q\|_{0,\Omega}^2\bigr)^{1/2}.
\]
引入双线性形式
\begin{equation}\label{eq:incompressible-sfvp-abstract-forms}
\left\{
\begin{aligned}
a((\mathbf u,\omega),(\mathbf v,\theta))
&=\nu(\omega,\theta),\\
b((\mathbf v,\theta),(\mu,q))
&=(\operatorname{curl}\mathbf v-\theta,\mu)
-(q,\operatorname{div}\mathbf v),
\end{aligned}
\right.
\end{equation}
以及右端项
\begin{equation}\label{eq:incompressible-sfvp-abstract-data}
\langle l,(\mathbf v,\theta)\rangle=(\mathbf f,\mathbf v),
\qquad
\chi=0.
\end{equation}

\MathBookSourcePage{214}
由这些数据得到的问题 $(Q)$ 为:

求 $(\mathbf u,\omega)\in X$ 和 $(\lambda,p)\in M$, 使
\begin{equation}\label{eq:incompressible-sfvp-four-field}
\left\{
\begin{aligned}
\nu(\omega,\theta)
+(\operatorname{curl}\mathbf v-\theta,\lambda)
-(p,\operatorname{div}\mathbf v)
&=(\mathbf f,\mathbf v)
&&\forall(\mathbf v,\theta)\in X,\\
(\operatorname{curl}\mathbf u-\omega,\mu)
-(q,\operatorname{div}\mathbf u)
&=0
&&\forall(\mu,q)\in M.
\end{aligned}
\right.
\end{equation}
其相联系的空间为
\[
V=\{(\mathbf v,\theta)\in X;\
\theta=\operatorname{curl}\mathbf v,\
\operatorname{div}\mathbf v=0\}.
\]
$a(\cdot,\cdot)$ 是 $V$-椭圆的,
而关于 $b(\cdot,\cdot)$ 的 inf-sup 条件由
Chapter~\ref{chap:stokes-foundations} 的结果直接得到.
因此有如下等价性结论.

\begin{Theorem}\label{thm:incompressible-sfvp-four-field-equivalence}
问题 \eqref{eq:incompressible-sfvp-four-field} 在
$X\times M$ 中有唯一解
$((\mathbf u,\omega),(\lambda,p))$.
此外, $(\mathbf u,p)$, 或等价地
$(\mathbf u,\omega,p)$, 是 Stokes 问题
\eqref{eq:incompressible-sfvp-stokes}, 或问题
\eqref{eq:incompressible-velocity-vorticity-pressure} 的解,
其中
\[
\omega=\operatorname{curl}\mathbf u,
\qquad
\lambda=\nu\omega.
\]
\end{Theorem}

\begin{Proof}
上述讨论蕴含问题
\eqref{eq:incompressible-sfvp-four-field} 有唯一解.
(唯一性也可由一个直接的简单论证得到.)
此外, 直接检验可知, 若
$(\mathbf u,\omega=\operatorname{curl}\mathbf u,p)$ 满足
\eqref{eq:incompressible-velocity-vorticity-pressure},
则 $((\mathbf u,\omega),(\nu\omega,p))$ 也满足
\eqref{eq:incompressible-sfvp-four-field}.
\end{Proof}

问题 \eqref{eq:incompressible-velocity-vorticity-pressure}
和 \eqref{eq:incompressible-sfvp-four-field}
也可用流函数表示.
回顾空间
\begin{equation}\label{eq:incompressible-stream-space-psi}
\Psi=\{\chi\in H^2(\Omega);\
\chi|_{\Gamma_0}=0,\
\chi|_{\Gamma_i}=c_i\text{ 为常数},\ 1\leq i\leq p,\
\partial\chi/\partial n|_\Gamma=0\}.
\end{equation}
由 Corollary~\ref{cor:sec3-v-stream-functions},
映射 $\operatorname{curl}$ 是从 $\Psi$ 到
\[
\{\mathbf v\in H_0^1(\Omega)^2;\
\operatorname{div}\mathbf v=0\}
\]
上的同构. 因而存在唯一的 $\psi\in\Psi$, 使
\[
\mathbf u=\operatorname{curl}\psi.
\]
在 \eqref{eq:incompressible-velocity-vorticity-pressure}
或 \eqref{eq:incompressible-sfvp-four-field} 中,
对所有 $\phi\in\Psi$ 取
$\mathbf v=\operatorname{curl}\phi$,
可知前一问题解的两个分量 $\mathbf u,\omega$,
或后一问题解的三个分量 $\mathbf u,\omega,\lambda$,
可分别由下列问题确定:

求 $(\psi,\omega)\in\Psi\times L^2(\Omega)$, 使
\begin{equation}\label{eq:incompressible-stream-vorticity-pair}
\left\{
\begin{aligned}
\nu(\omega,-\Delta\phi)
&=(\mathbf f,\operatorname{curl}\phi)
&&\forall\phi\in\Psi,\\
(\Delta\psi+\omega,\mu)&=0
&&\forall\mu\in L^2(\Omega);
\end{aligned}
\right.
\end{equation}
或者, 求
$(\psi,\omega,\lambda)\in
\Psi\times L^2(\Omega)\times L^2(\Omega)$, 使
\begin{equation}\label{eq:incompressible-stream-vorticity-multiplier}
\left\{
\begin{aligned}
\nu(\omega,\theta)-(\Delta\phi+\theta,\lambda)
&=(\mathbf f,\operatorname{curl}\phi)
&&\forall(\phi,\theta)\in\Psi\times L^2(\Omega),\\
(\Delta\psi+\omega,\mu)&=0
&&\forall\mu\in L^2(\Omega).
\end{aligned}
\right.
\end{equation}
这个问题属于 Subsection~\ref{subsec:abstract-general-result} 的框架,
其中
\begin{equation}\label{eq:incompressible-stream-abstract-interpretation}
\left\{
\begin{aligned}
X&=\Psi\times L^2(\Omega),
&M&=L^2(\Omega),\\
a((\psi,\omega),(\phi,\theta))
&=\nu(\omega,\theta),\\
b((\phi,\theta),\mu)
&=-(\Delta\phi+\theta,\mu),\\
\langle l,(\phi,\theta)\rangle
&=(\mathbf f,\operatorname{curl}\phi).
\end{aligned}
\right.
\end{equation}
\ChapterThreeSectionTwoDeferredContinuation
